\documentclass{article}
\usepackage[utf8]{inputenc}
\usepackage{graphicx}
\usepackage{url}
\usepackage[T1]{fontenc}
\usepackage{hyperref}
\usepackage{listings}
\usepackage{xcolor}
\usepackage{float}
\usepackage{algorithm}
\usepackage{algpseudocode}
\usepackage{amssymb}

\lstset{
  basicstyle=\ttfamily\footnotesize, 
  numbers=left,                     
  numberstyle=\tiny\color{gray},
  keywordstyle=\color{blue},        
  commentstyle=\color{green!50!black},
  stringstyle=\color{orange},        
  frame=single,                     
  breaklines=true,                   
  tabsize=2,
  showstringspaces=false,
  columns=flexible,
  frame=single,
  backgroundcolor=\color{gray!5}
}

\hypersetup{
    colorlinks=true,
    linkcolor=black,
    urlcolor=blue,
    citecolor=blue
}

\begin{document}

\begin{titlepage}
    \centering
    \vspace*{2cm}

    % Título principal
    {\Huge \bfseries Primera Entrega\par}
    \vspace{1.5cm}

    % Autores
    {\Large Gonzalo Pérez\par}
    \vspace{1cm}

    % Curso / asignatura
    {\large Ingeniería Informática\\
        Universidad del País Vasco (EHU) \\
        Desarrollo Avanzado de Software\par}
    \vspace{2cm}

    % Fecha
    {\large Marzo 2026\par}

    \vfill
\end{titlepage}


\tableofcontents
\newpage


\section{Introducción}

Como objetivo principal de esta primera entrega, se ha propuesto el desarrollo de una aplicación para el sistema operativo Android, utilizando para ello el entorno de desarrollo Android Studio\footnote{\url{https://developer.android.com/studio?hl=es-419}} y el lenguaje de programación Java\footnote{\url{https://www.java.com/es/}}. Los requisitos del proyecto establecen una serie de características obligatorias a implementar:

\begin{itemize}
    \item Uso de \textbf{ListView} o \textbf{RecyclerView} junto con \textbf{CardView} para mostrar conjuntos de elementos.
    \item Gestión de una base de datos local para listar, añadir y modificar elementos.
    \item Implementación de ventanas de diálogo.
    \item Generación de notificaciones locales.
    \item Gestión del ciclo de vida y control de la pila de actividades (\textit{Activities}).
\end{itemize}

Adicionalmente, se ha proporcionado una lista de características opcionales para enriquecer la aplicación:

\begin{itemize}
    \item Uso de \textbf{Fragments} para adaptar la interfaz visual según la orientación del dispositivo.
    \item Soporte multiidioma y opción para alternar entre diferentes idiomas.
    \item Lectura y escritura en ficheros de texto.
    \item Uso de preferencias (\textit{SharedPreferences}) para guardar de forma persistente ajustes del usuario.
    \item Definición de estilos y temas de personalización propios.
    \item Uso de \textbf{Intents} para interactuar con otras aplicaciones.
    \item Implementación de \textbf{Toolbar} y \textbf{Navigation Drawer} personalizados.
\end{itemize}

Para integrar todas estas funcionalidades y dar cumplimiento a los requisitos, se ha optado por diseñar y desarrollar una aplicación de gestión de tareas. Esta herramienta permite al usuario visualizar sus tareas, así como añadir nuevas, modificarlas, marcarlas como completadas o eliminarlas según sea necesario.

\section{Funcionalidades implementadas}

En este apartado se presentarán los elementos empleados y la función que cumplen dentro del proyecto.

\subsection{Elementos obligatorios}

\subsubsection{RecyclerView y CardView}

Se usa RecyclerView y CardView para mostrar los elementos de la aplicación. Para ello tambien se ha implementado un adaptador personalizado llamado \texttt{TareasAdapter}. Este adaptador se encarga de mostrar los elementos de la lista de tareas usando el layout \texttt{item\_tarea.xml}.

\subsubsection{Base de datos}

Se ha usando una base de datos local para almacenar las tareas. Para ello se ha utilziado SQLite. La base de datos se compone de una tabla llamada \texttt{tareas}, que tiene los siguientes campos:

\begin{itemize}
    \item id: identificador unico de la tarea
    \item titulo: titulo de la tarea
    \item descripcion: descripcion de la tarea
    \item fecha: fecha de la tarea
    \item prioridad: prioridad de la tarea donde: 0 = baja, 1 = media, 2 = alta
    \item completada: indica si la tarea esta completada donde 0 = no completada, 1 = completada
\end{itemize}

Este elemento es imprescindible ya que permite almacenar los datos de la aplicacion de forma persistente.

\subsubsection{Dialogos}

En la aplicación se utilizan tres tipos de dialogos:

\begin{itemize}
    \item Selector de fecha: Se utiliza para seleccionar la fecha de la tarea.
    \item Eliminar tareas: Para este proposito se han usado dos dialogos diferentes pero que cumplen la misma funcion. Uno para eliminar una tarea individual y otro para eliminar todas las tareas marcadas como completadas.
    \item Permiso de notificaciones: Este dialogo se utiliza cuando el usuario rechaza el permiso de notificaciones pero posteriormente intenta activarlas desde el menú de ajustes de la aplicación. En este caso se muestra un mensaje indicando que el permiso ha sido denegado y que debe ser activado desde los ajustes del sistema.
\end{itemize}

\subsubsection{Notificaciones}

En esta aplicacion las notificaiones cumplen la funcion de avisar al usuario de que tiene una tarea pendiente. Si los permisos de notificación están activados, se mostrará una notificación en la barra de notificaciones del dispositivo todos los dias a las 8 de la mañana con una lista de todas las tareas que tengan fecha de hoy o anterior. Si los permisos de notificación están desactivados, no se mostrará ninguna notificación.

\subsubsection{Control de la pila de actividades}

Este elemento es fundamental para el correcto funcionamiento de la aplicación, ya que permite gestionar el ciclo de vida de las actividades y controlar la pila de actividades. Para ello se ha asegurado de que la actividad principal sea la única actividad en la pila de actividades y que las actividades secundarias se muestren encima de la actividad principal, cerrandolas al volver a la actividad principal.

\newpage


\section{Clases y Base de datos}

\newpage

\section{Manual de Usuario}

\newpage

\section{Dificultades}

\newpage

\section{Fuentes utilizadas}

\newpage

\section{Repositorio de código}

\end{document}